%% algorithm.tex — Solution Methodology
%% Many-to-Many DRT Bidirectional Meeting Point Paper
%% Section 6 of the paper.

% ============================================================
\section{Solution Methodology}\label{sec:algorithm}
% ============================================================

The problem defined in Section~\ref{sec:model} is NP-hard (see
Section~\ref{subsec:complexity}). We develop two complementary solution approaches: an
MILP-based simplified ex-post diagnostic solved by Gurobi for small instances, and a rolling horizon
Adaptive Large Neighborhood Search (ALNS) heuristic for real-time operation at scale.

% ------------------------------------------------------------
\subsection{Simplified Ex-Post Diagnostic: MILP Formulation}\label{subsec:milp}
% ------------------------------------------------------------

For small instances, we formulate the static batch routing problem as a
Mixed-Integer Linear Program (MILP) to compute the optimal routing cost for a
given accepted set. The MILP serves as a \emph{simplified ex-post diagnostic over fixed accepted sets},
not a validation tool for the complete simulation framework, which includes
offer generation, routing, and stochastic passenger response.

\paragraph{Decision Variables.}
The MILP uses binary assignment variables $x_{r,m^P,m^D,v} \in \{0,1\}$, equal to 1 if
request $r$ is assigned to meeting point pair $(m^P, m^D) \in M_r^P \times M_r^D$ and
served by vehicle $v \in \mathcal{V}$. Continuous variables $s_{v,i} \geq 0$ represent
the scheduled service time at node $i$ in vehicle $v$'s route. The MILP is solved
only for the fixed realized set $\mathcal{R}^{\text{acc}}$; it does not optimize
the passenger response variable $z_r$.

\paragraph{Objective.}
The MILP minimizes the same weighted cost (Equation~\ref{eq:objective}), with each cost
component linearized in terms of $x_{r,m^P,m^D,v}$ and $s_{v,i}$.

\paragraph{Constraints.}
The formulation enforces the operational feasibility conditions
(\ref{con:capacity})--(\ref{con:domain-md}). Capacity constraint~(\ref{con:capacity}) is
linearized using standard load-tracking variables. Time-window
constraints~(\ref{con:tw-early})--(\ref{con:tw-late}), ride-time
constraint~(\ref{con:ridetime}), and route time consistency~(\ref{con:time-consistency})
are expressed as linear inequalities in $s_{v,i}$. Walking radius
constraints~(\ref{con:walk-pickup})--(\ref{con:walk-dropoff}) are enforced implicitly by
restricting the domain of $x_{r,m^P,m^D,v}$ to feasible meeting point pairs.

\paragraph{Solver and Scope.}
The MILP is solved using Gurobi~10 with a 300-second time limit. Any reported solver
gap is interpreted only within this simplified fixed-accepted-set diagnostic.

\textbf{Simplified ex-post diagnostic role.}
The MILP serves as a \emph{simplified ex-post diagnostic over fixed accepted sets},
not a real-time solver or a complete dynamic routing comparison.
Given the realized accepted set $\mathcal{R}^{\text{acc}}$ produced by a
completed ALNS run --- i.e., the set of requests for which the sampled
acceptance indicator $z_r = 1$ --- the MILP solves the optimal vehicle routing
for that \emph{fixed} accepted set.
Formally, the MILP treats $z_r$ as a known constant (equal to 1 for
$r \in \mathcal{R}^{\text{acc}}$ and 0 otherwise) and optimizes only over the
assignment variables $x_{r,m^P,m^D,v}$ and schedule variables $s_{v,i}$.
It does \emph{not} re-optimize the acceptance decisions: the stochastic
passenger responses are realized outcomes of the ALNS run, not MILP decision
variables.
The resulting MILP objective value is therefore a fixed-set diagnostic comparator, and the diagnostic gap
$\delta = (C^{\text{ALNS}} - C^{\text{MILP}}) / C^{\text{MILP}} \times 100\%$
describes the difference between the heuristic route and the simplified fixed-set
diagnostic solution only.
Instances with up to 30--50 passengers and 5--8 vehicles are tractable within
the time limit.

\begin{figure}[pos=h]
\centering
\includegraphics[width=\columnwidth]{figures/fig03_algorithm}
\caption{Online insertion and rolling horizon ALNS algorithm flowchart.}
\label{fig:algorithm}
\end{figure}

% ------------------------------------------------------------
\subsection{Heuristic Algorithm: Rolling Horizon ALNS}\label{subsec:alns}
% ------------------------------------------------------------

For real-time operation, we develop a rolling horizon ALNS heuristic that processes
arriving requests online and periodically re-optimizes committed routes. The algorithm
follows the three-layer architecture of Section~\ref{sec:three-layer}.

\paragraph{Online Insertion Mechanism.}
The ALNS heuristic uses a \emph{route-then-sample} mechanism that decouples
stochastic passenger acceptance from deterministic routing optimization.
At each request arrival, the system performs the following steps:
\begin{enumerate}
  \item \textbf{Bundle selection (Layer~1):} The system selects the best feasible bundle
        $b^*$ by minimizing the \emph{deterministic} incremental routing cost
        $\Delta C_{\text{routing}}(b, \{\sigma_v\})$, evaluated as the weighted sum
        $\alpha_1 \Delta C^{op} + \alpha_2 \Delta C^{wait} + \alpha_3 \Delta C^{walk}
        + \alpha_4 \Delta C^{IVT}$ of inserting bundle $b$ into the current vehicle
        schedules $\{\sigma_v\}$. The acceptance probability $P_{\text{accept}}(b^*)$
        is \emph{not} part of the bundle-selection objective.
  \item \textbf{Acceptance sampling (Layer~2):} After bundle selection, the binary logit
        acceptance probability $P_{\text{accept}}(b^*) = \exp(U_{rb^*}^k) /
        (1 + \exp(U_{rb^*}^k))$ (Equation~\ref{eq:binary-logit}) is computed, and
        the passenger's response is sampled as
        $z_r \sim \text{Bernoulli}(P_{\text{accept}}(b^*))$.
        If $z_r = 1$, the bundle is inserted into the vehicle route.
\end{enumerate}

\paragraph{Rationale for route-then-sample.}
This design separates the stochastic acceptance decision from the deterministic
cost minimization. An alternative approach that integrates $P_{\text{accept}}$
directly into the routing objective as $\mathbb{E}[\Delta C] = P_{\text{accept}}(b)
\times \Delta C_{\text{routing}}(b)$ would create a perverse incentive: bundles with
lower acceptance probability (worse service quality) would appear to have lower
expected routing cost, driving the system toward offering unattractive bundles.
The route-then-sample approach avoids this pathology while preserving the structural
benefit of passenger self-selection: passengers who would face high walking or
waiting disutility reject the service, and the system concentrates vehicle resources
on spatially compatible requests. Integrating $P_{\text{accept}}$ into the routing
objective in a welfare-consistent manner remains a direction for future work.

\paragraph{Rejection penalty (post-hoc only).}
The rejection penalty $\Gamma$ enters only the social welfare metric
$W = \sum_r [z_r U_{rb^*} - (1-z_r)\Gamma]$ (Section~\ref{sec:experiments}),
which is a post-hoc evaluation tool. $\Gamma$ does \emph{not} enter the bundle
selection, the routing objective, or the binary logit acceptance probability.
The acceptance probability $P_{\text{accept}}(b^*)$ is invariant to $\Gamma$.
This separation ensures that the operational decisions are consistent across
different post-hoc welfare-accounting settings encoded through $\Gamma$.

\subsubsection{Candidate Generation}

Upon arrival of request $r$, the system generates candidate meeting point sets $M_r^P$
and $M_r^D$ by filtering $\mathcal{M}$ against walking radii $\rho^P$ and $\rho^D$
(constraints~\ref{con:walk-pickup}--\ref{con:walk-dropoff}), then retaining the
top-$k^{\text{top}}$ candidates ranked by proximity to $o_r$ and $\delta_r$ respectively.
This yields $|M_r^P| \leq k^{\text{top}}$ pickup candidates and
$|M_r^D| \leq k^{\text{top}}$ dropoff candidates. The filtering step runs in
$O(|\mathcal{M}| \log |\mathcal{M}|)$ per request using a pre-sorted distance index.

\subsubsection{Online Insertion Evaluation}

For each bundle $(m^P, m^D) \in M_r^P \times M_r^D$, the algorithm enumerates all
vehicles $v \in \mathcal{V}$ and all insertion position pairs $(\pi^P, \pi^D)$ with
$\pi^D > \pi^P$ (constraint~\ref{con:precedence-pos}). Feasibility against
constraints~(\ref{con:capacity}), (\ref{con:tw-early})--(\ref{con:tw-late}),
(\ref{con:ridetime}), and (\ref{con:time-consistency}) is checked in $O(1)$ per position
pair using forward and backward schedule arrays maintained for each vehicle route. The
best feasible insertion is selected by minimum incremental cost $\Delta C$ to the
objective~(\ref{eq:objective}). After selecting $b^*$, the binary logit acceptance probability
$P_{\text{accept}}(b^*) = \exp(U_{rb^*}^k) / (1 + \exp(U_{rb^*}^k))$
(Equation~\ref{eq:binary-logit}) is computed and the passenger's response is sampled
as $z_r \sim \text{Bernoulli}(P_{\text{accept}}(b^*))$.

\subsubsection{ALNS Destroy and Repair Operators}

Every $\Delta$ minutes, the rolling horizon re-optimization applies ALNS
\citep{ropke2006} to the active request set $\mathcal{R}^{\text{act}}(t^{\text{now}})$.
The algorithm uses five operators:

\begin{enumerate}
  \item \textbf{Request reassignment:} Remove $k$ randomly selected requests from their
        current vehicle routes and reinsert them greedily into the best feasible positions
        across all vehicles.
  \item \textbf{Pickup/dropoff pair swap:} Exchange the meeting point pair
        $(m_r^P, m_r^D)$ assignments between two requests $r$ and $r'$, updating
        insertion positions accordingly.
  \item \textbf{Insertion position exchange:} Swap the route positions of two service
        nodes on the same vehicle, subject to precedence and feasibility.
  \item \textbf{Route segment reconstruction:} Remove a contiguous route segment from a
        vehicle, then reinsert the affected requests optimally across all vehicles.
  \item \textbf{Worst removal:} Remove the $k$ requests with the highest marginal cost
        contribution to the objective, then reinsert greedily.
\end{enumerate}

Operator selection follows the standard ALNS roulette-wheel mechanism
\citep{ropke2006}: each operator maintains a score updated by its historical improvement
contribution, and operators are selected with probability proportional to their score.
Meeting point reassignment is integrated into the repair step of operators 1, 4, and 5:
when a request is reinserted, the algorithm re-evaluates all $(m^P, m^D)$ pairs in
$M_r^P \times M_r^D$ and selects the pair that minimizes insertion cost.

\subsubsection{Rolling Horizon Loop}

The rolling horizon framework \citep{bent2004} triggers ALNS re-optimization at times
$t^{\text{now}} = 0, \Delta, 2\Delta, \ldots$ (Equation~\ref{eq:rh-trigger}). At each
trigger, the algorithm operates on the active set
$\mathcal{R}^{\text{act}}(t^{\text{now}})$ over the window
$[t^{\text{now}},\; t^{\text{now}} + H]$ with $H = 30$ minutes and $\Delta = 5$ minutes.
Committed nodes (already served or within the commitment horizon) are locked and cannot
be reassigned.

\paragraph{Commitment assumption.}
Once a service offer bundle $b_r^* = (m_r^P, m_r^D, \tau_r, p_r)$ is communicated to
passenger $r$ and accepted ($z_r = 1$), the assigned meeting points and scheduled pickup
time are \emph{committed}: they are not subject to re-optimization in subsequent rolling
horizon cycles. This commitment is operationally necessary --- passengers who have begun
walking to a meeting point cannot be redirected --- and is enforced by locking the
corresponding nodes in the route graph before each ALNS re-optimization trigger.
Formally, a node $i \in \sigma_v$ is committed if its scheduled service time satisfies
$\tau_i \leq t^{\text{now}} + \delta_{\text{commit}}$, where $\delta_{\text{commit}} = 5$
minutes is the commitment horizon (equal to the re-optimization interval $\Delta$).
Uncommitted accepted nodes may be re-sequenced within the vehicle route to improve
efficiency, but their assigned meeting points remain fixed.

The ALNS runs for 50 iterations per re-optimization trigger, with each
iteration applying one destroy--repair operator pair. After re-optimization, updated
routes are committed for the next $\Delta$ minutes and fed back to Layer~1 for subsequent
request evaluations.

\begin{algorithm}[h]
\caption{Rolling Horizon ALNS for Bidirectional DARPmp}
\label{alg:rh-alns}
\begin{algorithmic}[1]
\Require Request stream $\mathcal{R}$, fleet $\mathcal{V}$, meeting points $\mathcal{M}$,
         horizon $H$, re-opt interval $\Delta$, rejection penalty $\Gamma$
\Ensure  Accepted set $\mathcal{R}^{\text{acc}}$, routes $\{\sigma_v\}_{v \in \mathcal{V}}$

\State Initialize routes $\sigma_v \leftarrow \emptyset$ for all $v \in \mathcal{V}$
\State $t^{\text{now}} \leftarrow 0$

\While{request stream not exhausted}
  \State \textbf{// --- Online dispatch: process next arriving request ---}
  \State Receive request $r$ at time $t_r^{\text{req}}$
  \State $M_r^P \leftarrow \{m \in \mathcal{M} \mid d(o_r, m) \leq \rho^P\}$,
         $M_r^D \leftarrow \{m \in \mathcal{M} \mid d(\delta_r, m) \leq \rho^D\}$
         \Comment{Candidate meeting points}
  \State $b^* \leftarrow \textsc{BestInsertion}(r, M_r^P, M_r^D, \{\sigma_v\})$
         \Comment{Layer 1: select best feasible bundle}
  \If{$b^* = \emptyset$}
    \State $z_r \leftarrow 0$; record rejection
    \Comment{No feasible insertion found}
  \Else
    \State $p_{\text{acc}} \leftarrow \dfrac{\exp(U_{rb^*}^{k_r})}{1 + \exp(U_{rb^*}^{k_r})}$
           \Comment{Binary logit acceptance (Eq.~\ref{eq:binary-logit})}
    \State $z_r \sim \text{Bernoulli}(p_{\text{acc}})$
           \Comment{Layer 2: passenger response}
    \If{$z_r = 1$}
      \State Insert $b^*$ into $\sigma_{v_{b^*}}$ at positions $(\pi^P, \pi^D)$
      \State $\mathcal{R}^{\text{acc}} \leftarrow \mathcal{R}^{\text{acc}} \cup \{r\}$
    \EndIf
  \EndIf

  \State \textbf{// --- Periodic re-optimization ---}
  \If{$t^{\text{now}} \bmod \Delta = 0$}
    \State $\mathcal{R}^{\text{act}} \leftarrow \{r' \in \mathcal{R}^{\text{acc}} \mid r' \text{ not yet completed}\}$
    \State $\{\sigma_v\} \leftarrow \textsc{ALNS}(\mathcal{R}^{\text{act}}, \{\sigma_v\}, H)$
           \Comment{Layer 3: rolling horizon re-opt}
  \EndIf
  \State $t^{\text{now}} \leftarrow t^{\text{now}} + \delta t$
\EndWhile
\State \Return $\mathcal{R}^{\text{acc}},\; \{\sigma_v\}$
\end{algorithmic}
\end{algorithm}

% ------------------------------------------------------------
\subsection{Computational Complexity}\label{subsec:complexity}
% ------------------------------------------------------------

\paragraph{NP-Hardness.}
The many-to-many DRT problem with bidirectional meeting points is NP-hard. This follows
by reduction from the Dial-a-Ride Problem (DARP) \citep{cordeau2007}: for each request
$r$, set the candidate pickup set to a singleton containing only the passenger's origin
($M_r^P = \{o_r\}$) and the candidate dropoff set to a singleton containing only the
destination ($M_r^D = \{\delta_r\}$), with $\rho^P$ and $\rho^D$ set sufficiently large
to include these singletons. The problem then reduces to standard DARP (door-to-door
with no meeting point choice), which is NP-hard. The bidirectional meeting-point
extension strictly generalizes this case by allowing $|M_r^P|, |M_r^D| \geq 1$, adding
the joint optimization of $m_r^P$ and $m_r^D$, so the problem is at least as hard.

\paragraph{Heuristic Complexity.}
The per-request candidate generation step runs in $O(|\mathcal{M}| \log |\mathcal{M}|)$.
Online insertion evaluation for a single request costs
$O(|\mathcal{V}| \cdot |\sigma_v|^2 \cdot (k^{\text{top}})^2)$, where $|\sigma_v|$ is
the average route length. Each ALNS iteration costs
$O(n \cdot |\mathcal{V}| \cdot |\sigma_v|)$ where $n = |\mathcal{R}^{\text{act}}|$.
The rolling horizon bounds the active set size: at any trigger, at most
$H / \bar{\lambda}^{-1}$ requests are active, where $\bar{\lambda}$ is the average
arrival rate. This keeps per-step complexity manageable even as the total number of
requests grows.

\paragraph{Practical Performance.}
On the 200-request, 15-vehicle simulation instances, the average decision time per
request is below 1 second, satisfying the real-time requirement of online DRT operation.
The simplified ex-post diagnostic over fixed accepted sets (Section~\ref{subsec:milp})
compares ALNS routing with the diagnostic MILP solution for a fixed accepted set on small
instances. These diagnostic differences reflect the limited fixed-set setup, not an
assessment of overall solution quality, since the MILP does not model stochastic
passenger response.
